\clearpage
\section{Introduction}\label{sec:intro}
\par The conduct of monetary policy by central banks relies fundamentally on the ability to forecast key macroeconomic variables and communicate policy intentions to market participants. In the European context, the European Central Bank's (ECB) deposit facility rate serves as the primary instrument for steering monetary policy in the euro area. Accurate forecasting of policy rates is crucial not only for central bank decision-making but also for financial market participants, governments, and businesses that must anticipate future monetary conditions to make informed decisions.
\par This paper develops and evaluates a forecasting methodology for the ECB deposit facility rate based on the Taylor rule framework. The Taylor rule, originally proposed by \cite{taylor1993discretion}, provides a systematic approach to understanding central bank behaviour by relating the policy interest rate to deviations of inflation from target and output from potential. Our methodology combines the theoretical foundation of the Taylor rule with time series forecasting techniques. Specifically, we employ ARIMA models to generate forecasts of the Taylor rule inputs—the inflation gap (deviation of HICP inflation from the ECB's 2\% target) and the output gap (estimated using the Hamilton filter applied to euro area real GDP), which are then combined to predict future policy rates. While we estimate specifications using both actual inflation and inflation expectations from the ECB Survey of Professional Forecasters, our primary results focus on models using actual inflation data. The forecasting challenge is particularly acute in the post-2008 financial crisis environment, which has been characterized by the zero lower bound (ZLB) constraint, unconventional monetary policies, and significant structural breaks. To address these challenges, we incorporate the Wu-Xia shadow rate \citep{wu2016measuring}, which captures the stance of monetary policy when conventional rates are constrained by the ZLB. We also employ rolling-window estimation techniques with a window size of 85 quarters (approximately 21 years) to account for time-varying relationships between policy rates and macroeconomic fundamentals.
\par We evaluate our forecasting methodology using quarterly data from 1999 Q1 to 2025 Q3 through a rigorous pseudo out-of-sample framework, conducting forecasts up to ten quarters ahead using a rolling window approach. The evaluation employs both absolute performance metrics (Mincer-Zarnowitz tests for bias and efficiency) and relative performance metrics comparing mean squared forecast errors against a benchmark ARIMA model using Diebold-Mariano tests. This comprehensive evaluation framework allows us to assess the practical utility of Taylor rule-based forecasts relative to atheoretical time series benchmarks. Our analysis reveals several key findings. First, we document significant structural breaks in the Taylor rule relationship, particularly around the onset of the ZLB period (2008 Q3) and the COVID-19 pandemic (2020 Q1), which motivate our use of rolling-window estimation. Second, we find that Taylor rule-based forecasts can provide competitive forecasting performance relative to benchmark ARIMA models, particularly at medium-to-long horizons. Third, we demonstrate that incorporating interest rate persistence through lagged dependent variables substantially improves forecast accuracy. Finally, we provide actual forecasts for the ECB deposit facility rate extending to 2028 Q1, accompanied by prediction intervals with uncertainty derived from the pseudo out-of-sample evaluation exercise.

\section{Literature Review}\label{sec:litreview}
\par The Taylor rule has become one of the most influential frameworks in monetary economics since its introduction by \cite{taylor1993discretion}. The original formulation proposed that the Federal Reserve's policy rate could be well approximated by a simple linear function of the inflation gap the deviation of actual inflation from the central bank's target and the output gap the percentage deviation of real GDP from potential GDP. The elegance and empirical success of this simple rule sparked an extensive literature examining its applicability across different countries, time periods, and policy regimes. \cite{clarida2000monetary} extended the Taylor rule framework by estimating forward-looking specifications for the Federal Reserve, Bundesbank, and Bank of Japan, demonstrating that central banks respond to expected rather than contemporaneous inflation. \cite{orphanides2001monetary} highlighted the importance of real-time data, showing that Taylor rule estimates can be substantially affected by data revisions and the mismeasurement of potential output. For the ECB specifically, \cite{gerdesmeier2004taylor} and \cite{fendel2009inflation} found that simple Taylor rules provide a reasonable description of ECB policy behaviour during the pre-crisis period, though with coefficients differing from Taylor's original calibration. \cite{gorter2008taylor} provided important evidence supporting the forward-looking nature of ECB policy, demonstrating that Taylor rules estimated using expectations data from Consensus Economics for inflation and output growth significantly outperform specifications based on realized outcomes, with the ECB taking expected inflation and expected output growth into account when setting interest rates. \cite{klose2023estimated} developed a comprehensive real-time database for the ECB decision dates with daily forecasts for inflation and output up to fourteen months ahead, finding that the ECB's inflation forecast horizon is rather long-term (over one year) while the picture for output is less clear, and demonstrating that these forecast horizons tend to be time-varying across different policy regimes.
\par The global financial crisis and subsequent period of extremely low interest rates posed significant challenges for both the implementation and estimation of Taylor rules. When policy rates are constrained by the ZLB, conventional Taylor rule estimates may fail to capture the true stance of monetary policy. \cite{wu2016measuring} addressed this issue by developing shadow rate models that estimate a hypothetical policy rate in the absence of the ZLB constraint. Their methodology, which we employ in this paper, provides a continuous measure of monetary policy stance throughout both conventional and unconventional policy regimes. \cite{krippner2013measuring} proposed an alternative framework based on a ``shadow short rate'' that uses a tractable approximation to model the term structure in ZLB environments, demonstrating its application to Japanese monetary policy and showing that negative shadow rates can effectively indicate the stance of unconventional monetary policy. \cite{lombardi2018taylor} developed an alternative shadow policy rate for the United States that facilitates assessment of monetary policy stance against Taylor rule benchmarks during the ZLB period, providing a comprehensive measure of Federal Reserve policy actions including unconventional measures. \cite{hofmann2012taylor} documented that policy rates in both advanced and emerging market economies had been systematically below levels implied by the Taylor rule since the early 2000s, a pattern they termed the global ``Great Deviation'', reflecting either accommodative monetary policy or potentially lower equilibrium real interest rates.
\par A substantial body of research has documented that the parameters of the Taylor rule are not constant over time but exhibit significant variation across policy regimes. \cite{kim2006estimation} estimated a forward-looking monetary policy rule with time-varying parameters for the Federal Reserve, demonstrating that policy responsiveness to inflation and output has changed substantially over time. \cite{trecroci2010monetary} estimated time-varying interest rate rules for major OECD economies including the euro area, finding that monetary policy regime shifts are pervasive and that time-varying parameter specifications substantially outperform constant-parameter Taylor rules in tracking actual policy rates. \cite{sturm2011communication} examined whether ECB communication adds information beyond that contained in a forward-looking Taylor rule model with expected inflation and output growth, finding that communication indicators based on the ECB President's statements provide significant additional predictive power for forecasting interest rate decisions, even after controlling for market expectations embedded in interbank rates. \cite{czudaj2023forecasts} examined whether professional forecasters form their expectations regarding the ECB policy rate consistent with the Taylor rule, using micro-level data from the ECB Survey of Professional Forecasters, finding that professionals indeed form expectations in line with the Taylor rule, though this connection has diminished over time, especially after the policy rate hit the zero lower bound.
\par The use of policy rules for interest rate forecasting has received considerable attention. \cite{molodtsova2009taylor} examined the out-of-sample forecasting performance of Taylor rule models for exchange rates, finding that Taylor rule fundamentals could outperform random walk benchmarks. \cite{rossi2013advances} provided a comprehensive review of exchange rate forecasting methods, highlighting that the relative performance of different approaches can be highly sensitive to the evaluation period and forecast horizon. For direct interest rate forecasting, \cite{ang2003noarbitrage} compared various approaches including affine term structure models, macro-finance models, and time series benchmarks, generally finding that simple time series models were difficult to outperform at short horizons. However, \cite{mccracken2016fred} demonstrated that combining forecasts from multiple models, including those based on policy rules, could improve forecast accuracy. The two-step forecasting approach we employ—forecasting the inputs to the Taylor rule and then combining them to generate interest rate forecasts—relates to the literature on forecasting with generated regressors. \cite{pagan1984econometric} demonstrated that using generated regressors in two-step estimation procedures requires careful treatment to ensure valid statistical inference and can introduce additional forecast error variance.
\par Our approach addresses both parameter instability and forecast evaluation rigorously. We employ rolling-window estimation to allow Taylor rule coefficients to evolve gradually over time, complemented by formal structural break tests \citep{chowTestsEqualitySets1960, baiEstimatingTestingLinear1998} to verify the presence and timing of regime changes. Our evaluation framework employs both Mincer-Zarnowitz tests for absolute performance and Diebold-Mariano tests for relative performance, drawing on the forecast evaluation literature developed by \cite{diebold1995comparing} and \cite{west1996asymptotic}, allowing us to assess empirically whether the structure imposed by the Taylor rule framework provides sufficient benefits to offset any additional forecast error variance from the two-step procedure.
